% !TeX encoding = UTF-8
% Short abstract for the thesis. Edit this file with the final abstract text.
\noindent

Large Language Models (LLMs) are increasingly deployed in healthcare applications; however, their reliability and safety in multi-turn medical dialogues remain insufficiently understood. 
Hallucination, contextual degradation, and the omission of critical information pose significant risks in clinical settings, where existing single-turn evaluation metrics fail to capture the full complexity of chatbot performance.

This thesis proposes and evaluates a comprehensive benchmarking pipeline for LLM-based medical chatbots, developed in the context of the Patient Information Assistant (PIA) project for prenatal care consultations. 
The framework addresses four evaluation targets: information extraction quality, topic coverage, citation faithfulness, and mandatory question compliance. 
For information extraction, four prompting strategies were benchmarked across four language models using thresholded semantic matching and the adapted SUSWIR metric, identifying GPT-5 Mini with Naive prompting as the optimal configuration. 
Topic coverage was evaluated using three complementary methods: Hungarian matching, bi-encoder with entailment, and LLM-as-a-Judge, across two synthetic conversation datasets of 60 and 216 conversations respectively. 
Citation faithfulness was assessed using a two-level strict and relaxed precision framework. 
Finally, mandatory question compliance was evaluated by comparing naive and supervised dialogue management configurations.

The results demonstrate that supervised dialogue management substantially improves topic coverage and raises mandatory question compliance to 100\% in the smaller dataset, while increasing question recall by 21 percentage points in the larger dataset, without degrading citation faithfulness. 
These findings provide empirical support for the necessity of active conversational scaffolding in safety-critical medical applications and offer a modular, extensible evaluation blueprint adaptable to future clinical domains and patient populations.


